\section{Diseño Metodológico}

El diseño metodológico se organiza a partir de los objetivos específicos, relacionando actividades, recursos, productos esperados y tiempo estimado. El proyecto corresponde a una investigación aplicada, ya que busca desarrollar un prototipo tecnológico orientado al aprendizaje práctico de ciberseguridad.

\begin{longtable}{|p{0.23\textwidth}|p{0.25\textwidth}|p{0.18\textwidth}|p{0.20\textwidth}|p{0.08\textwidth}|}
\hline
\textbf{Objetivo específico} & \textbf{Actividades} & \textbf{Recursos} & \textbf{Producto esperado} & \textbf{Tiempo} \\
\hline
Diseñar escenarios interactivos que simulen situaciones de ataque cibernético. &
Definir casos de \textit{phishing}, malware e ingeniería social. Estructurar niveles y situaciones de decisión. &
Bibliografía, casos de amenaza, guías de ciberseguridad. &
Documento de escenarios y flujo de niveles. &
3 sem. \\
\hline
Incorporar mecánicas de juego basadas en toma de decisiones. &
Diseñar acciones del jugador, reglas de decisión, estados de riesgo y consecuencias. &
Referentes de juegos serios, prototipos de interfaz, motor de desarrollo. &
Sistema inicial de mecánicas jugables. &
3 sem. \\
\hline
Organizar conceptos de ciberseguridad en actividades progresivas. &
Ordenar contenidos por dificultad. Relacionar amenazas, vulnerabilidades y medidas de protección. &
Marco teórico, contenidos de seguridad, criterios pedagógicos. &
Ruta progresiva de aprendizaje dentro del videojuego. &
2 sem. \\
\hline
Relacionar situaciones reales de riesgo digital con ejercicios prácticos. &
Construir casos representativos con correos, alertas, llamadas y eventos simulados. &
Ejemplos de incidentes, patrones de ataque, documentación técnica. &
Banco de casos prácticos para el prototipo. &
3 sem. \\
\hline
Implementar retroalimentación al finalizar cada escenario. &
Diseñar reportes de desempeño, consecuencias narrativas y resumen de decisiones. &
Registros del juego, criterios de evaluación, diseño de interfaz. &
Módulo de retroalimentación por escenario. &
3 sem. \\
\hline
Realizar pruebas de usabilidad con usuarios. &
Preparar guía de prueba, aplicar evaluación con usuarios y analizar observaciones. &
Usuarios de prueba, cuestionario, prototipo funcional. &
Reporte de usabilidad y ajustes del prototipo. &
3 sem. \\
\hline
\end{longtable}
